%%%%%%%%%%%%%%%%%%%%%%%%%%%%%%%%%
%
%       EXERCÍCIO
%
%%%%%%%%%%%%%%%%%%%%%%%%%%%%%%%%%

\ifdefstring{\atividade}{prova}{%
    \ifdefstring{\modo}{objetivo}{%
        \renewcommand{\valorquestao}{\ValorQObj\ ponto}
    }{%
        \renewcommand{\valorquestao}{\ValorQDisc\ pontos}
    }%
}%

\begin{exercicioBanco}[\valorquestao]
Para cada afirmativa abaixo, assinale se ela é verdadeira (V) ou falsa (F).
%
\begin{enumerate}[\(\bullet\)]
    \item \(A = (1,2)\), \(B = (3,6)\) e \(C = (5,10)\) estão alinhados.
    %
    \item As retas \(r\) e \(s\) determinadas, respectivamente, pelas equações \(6x + 9y - 3 = 0\) e \(2x + 3y + 5 = 0\) são paralelas.
    %
    \item A distância entre o ponto \(P = (1,2)\) e a reta \(r: 3x + 4y - 10 = 0\) é \(1\).
\end{enumerate}
%

% Define as alternativas
\newcommand{\alternativas}{%
    \begin{center}
        \begin{tabularx}{\textwidth}{XXXXX}
            (a) (V, V, V). &
            (b) (V, V, F). &
            (c) (F, V, F). &
            (d) (V, F, V). &
            (e) (F, V, V).
        \end{tabularx}
    \end{center}
}

% Define a resposta correta
\newcommand{\resposta}{B}

% Lógica condicional para exibição
\ifdefstring{\atividade}{lista}{%
    \alternativas
    \vspace{0.5em}
    
    \noindent\textbf{Resposta:} letra \textbf{\resposta}.
}{%
    \ifdefstring{\modo}{objetivo}{%
        \alternativas
    }{%
        % modo = discursiva → não mostra alternativas
    }
}
\end{exercicioBanco}

